\chapter{Readiness solutions}

\label{learn:solutions}

Compare the operations as well as the final value. Return to the named lesson when a step needs repair.

\section{Solutions: From convincing examples to proof}

\label{solution:entry:proof}

\textbf{Entry answer.} Distribute multiplication to obtain x\textasciicircum{}2+2x+1.\par

\label{solution:ready:proof}

\textbf{Readiness answer.} Write the integers as 2m+1 and 2n+1 with integer m,n. Their sum is 2(m+n+1), twice an integer.\par

\hyperref[learn:proof]{Return to the lesson} or \hyperref[readiness:proof]{return to its next-question check}.

\section{Solutions: From balanced equations to matrix maps}

\label{solution:entry:matrix_maps}

\textbf{Entry answer.} Subtract the first equation from the second: x=4. Substitute to get y=3.\par

\label{solution:ready:matrix_maps}

\textbf{Readiness answer.} The outputs are 11 and 12. Their difference is -1, also obtained from the composite row (2,-1) applied to (2,5).\par

\hyperref[learn:matrix_maps]{Return to the lesson} or \hyperref[readiness:matrix_maps]{return to its next-question check}.

\section{Solutions: Complex numbers: multiplication, length, and conjugation}

\label{solution:entry:complex_numbers}

\textbf{Entry answer.} Distribution gives ac+ad+bc+bd; complex multiplication also uses i\textasciicircum{}2=-1.\par

\label{solution:ready:complex_numbers}

\textbf{Readiness answer.} The conjugate is 3-4i, squared modulus is 25, and reciprocal is (3-4i)/25. Multiplication checks a product of one.\par

\hyperref[learn:complex_numbers]{Return to the lesson} or \hyperref[readiness:complex_numbers]{return to its next-question check}.

\section{Solutions: Linear algebra: coordinates, maps, and modes}

\label{solution:entry:linear_algebra}

\textbf{Entry answer.} The output is (1,-1).\par

\label{solution:ready:linear_algebra}

\textbf{Readiness answer.} A(1,1)=3(1,1) and A(1,-1)=(1,-1). The two vectors are independent because neither is a scalar multiple of the other.\par

\hyperref[learn:linear_algebra]{Return to the lesson} or \hyperref[readiness:linear_algebra]{return to its next-question check}.

\section{Solutions: Calculus: local change and accumulated change}

\label{solution:entry:calculus}

\textbf{Entry answer.} The geometric tail is 2\textasciicircum{}(1-n), tending to zero as n increases.\par

\label{solution:ready:calculus}

\textbf{Readiness answer.} Subtract x\textasciicircum{}2 from (x+h)\textasciicircum{}2 and divide by h to obtain 2x+h. Let h approach zero to obtain 2x.\par

\hyperref[learn:calculus]{Return to the lesson} or \hyperref[readiness:calculus]{return to its next-question check}.

\section{Solutions: Several variables: local maps and conservation}

\label{solution:entry:multivariable}

\textbf{Entry answer.} The derivative is 2x; the dot product is 11.\par

\label{solution:ready:multivariable}

\textbf{Readiness answer.} The partial derivatives are 2x and 6y, giving gradient (2,12). Its dot product with a unit direction gives the directional rate.\par

\hyperref[learn:multivariable]{Return to the lesson} or \hyperref[readiness:multivariable]{return to its next-question check}.

\section{Solutions: Dynamics: equations, modes, and trustworthy computation}

\label{solution:entry:dynamics}

\textbf{Entry answer.} The chain rule gives -2 exp(-2t).\par

\label{solution:ready:dynamics}

\textbf{Readiness answer.} Solve -1\textless{}1-2h\textless{}1 to obtain 0\textless{}h\textless{}1. At h=1 the factor is -1 and the numerical solution does not decay.\par

\hyperref[learn:dynamics]{Return to the lesson} or \hyperref[readiness:dynamics]{return to its next-question check}.

\section{Solutions: Probability: uncertainty, evidence, and comparison}

\label{solution:entry:probability}

\textbf{Entry answer.} The fraction is 4/100=1/25.\par

\label{solution:ready:probability}

\textbf{Readiness answer.} There are 90+495=585 positives. The posterior is 90/585=2/13. The denominator counts observed positives, rather than the whole population.\par

\hyperref[learn:probability]{Return to the lesson} or \hyperref[readiness:probability]{return to its next-question check}.

\section{Solutions: Numerical reasoning: approximation with an error budget}

\label{solution:entry:numerical}

\textbf{Entry answer.} The discrete factor changes with h; its repeated product approximates the exponential only in a controlled limit.\par

\label{solution:ready:numerical}

\textbf{Readiness answer.} One panel gives 1/2 with error 1/6. Two give 3/8 with error 1/24, a factor-four reduction for this exact example.\par

\hyperref[learn:numerical]{Return to the lesson} or \hyperref[readiness:numerical]{return to its next-question check}.

\section{Solutions: Algebraic structure: symmetry, brackets, and roots}

\label{solution:entry:algebra}

\textbf{Entry answer.} The first basis vector selects the first column and the second selects the second.\par

\label{solution:ready:algebra}

\textbf{Readiness answer.} AB is diag(1,0), BA is diag(0,1), so the commutator is diag(1,-1). The nonzero result records order dependence.\par

\hyperref[learn:algebra]{Return to the lesson} or \hyperref[readiness:algebra]{return to its next-question check}.

\section{Solutions: Geometry and analysis: distance, dimension, and complex structure}

\label{solution:entry:geometry}

\textbf{Entry answer.} The modulus is 5, obtained from the square root of 3\textasciicircum{}2+4\textasciicircum{}2.\par

\label{solution:ready:geometry}

\textbf{Readiness answer.} There are 2\textasciicircum{}n intervals of length 3\textasciicircum{}(-n). The self-similar dimension solves 2*(1/3)\textasciicircum{}d=1, so d=log(2)/log(3). The chapter explains the additional dimension argument.\par

\hyperref[learn:geometry]{Return to the lesson} or \hyperref[readiness:geometry]{return to its next-question check}.

\section{Solutions: Doubling changes multiplication as well as dimension}

\label{solution:entry:hypercomplex}

\textbf{Entry answer.} The cross terms cancel and i\textasciicircum{}2=-1, giving 2.\par

\label{solution:ready:hypercomplex}

\textbf{Readiness answer.} In a division algebra, left multiplication by any nonzero element is invertible and hence injective. A nonzero element mapped to zero violates injectivity.\par

\hyperref[learn:hypercomplex]{Return to the lesson} or \hyperref[readiness:hypercomplex]{return to its next-question check}.

\section{Solutions: Worked example: probabilities from a complex state}

\label{solution:entry:quantum}

\textbf{Entry answer.} Conjugate-first multiplication gives 1+1=2.\par

\label{solution:ready:quantum}

\textbf{Readiness answer.} Divide by sqrt(2). Each squared modulus is 1/2, summing to one. Interpreting those numbers as measurement probabilities uses the Born postulate.\par

\hyperref[learn:quantum]{Return to the lesson} or \hyperref[readiness:quantum]{return to its next-question check}.

\section{Solutions: A lattice-Boltzmann step has exact moment obligations}

\label{solution:entry:fluids}

\textbf{Entry answer.} Roundoff comes from finite arithmetic; truncation comes from replacing a limiting operation by a finite approximation.\par

\label{solution:ready:fluids}

\textbf{Readiness answer.} The sum is 4/9+4/9+1/9=1. Opposite velocities have equal weights, so each component of the first moment cancels.\par

\hyperref[learn:fluids]{Return to the lesson} or \hyperref[readiness:fluids]{return to its next-question check}.
